% Part: turing-machines
% Chapter: machines-computations
% Section: representing-tms

\documentclass[../../../include/open-logic-section]{subfiles}

\begin{document}

\olfileid[fr]{tur}{mac}{rep}
\olsection{Représenter les machines de Turing}

\begin{explain}
On peut représenter visuellement une machine de
Turing par un \emph{diagramme d'états}. Celui-ci
est formé de nœuds reliés par des flèches.
Chaque nœud représente un état de la machine;
chaque flèche, une instruction exécutable à partir
de cet état. Le détail de l'instruction figure
au-dessus ou au-dessous de la flèche correspondante.
Considérons la machine suivante, qui n'a que deux
états internes, $q_0$ et~$q_1$, et une seule
instruction :
\[
\begin{tikzpicture}[->,>=stealth',shorten >=1pt,auto,node distance=2.8cm,
                    semithick]
  \tikzstyle{every state}=[fill=none,draw=black,text=black]

  \node[initial,state]   (A)              {$q_0$};
  \node[state]   (B) [right of=A] {$q_1$};

  \path (A) edge  node {\TMtrans{\TMblank}{\TMstroke}{\TMright}} (B);
\end{tikzpicture}
\]
Rappelons que la machine possède une tête de
lecture-écriture et un ruban portant l'entrée.
L'instruction se lit ainsi : \emph{si, dans
l'état~$q_0$, la tête lit~$\TMblank$, écrire
$\TMstroke$, déplacer la tête vers la droite
et passer dans l'état~$q_1$}. Autrement dit,
la fonction de transition associe à
$\tuple{q_0,\TMblank}$ le triplet
$\tuple{q_1,\TMstroke,\TMright}$.
\end{explain}

\begin{ex}
\emph{Machine paire} : la machine de Turing suivante
s'arrête si et seulement si le ruban porte un nombre
pair de symboles $\TMstroke$, à condition que tous
ces symboles précèdent le premier $\TMblank$ du
ruban.
\[
\begin{tikzpicture}[->,>=stealth',shorten >=1pt,auto,node distance=2.8cm,
                    semithick]
  \tikzstyle{every state}=[fill=none,draw=black,text=black]

  \node[initial,state]         (A)              {$q_0$};
  \node[state]         (B) [right of=A] {$q_1$};

  \path (A) edge [bend left] node {\TMtrans{\TMstroke}{\TMstroke}{\TMright}} (B)
        (B) edge [loop above] node {\TMtrans{\TMblank}{\TMblank}{\TMright}} (B)
            edge [bend left] node {\TMtrans{\TMstroke}{\TMstroke}{\TMright}} (A);
\end{tikzpicture}
\]

Le diagramme correspond à la fonction de transition
suivante :
\begin{align*}
\delta(q_0, \TMstroke) & = \tuple{q_1, \TMstroke, \TMright},\\
\delta(q_1, \TMstroke) & = \tuple{q_0, \TMstroke, \TMright},\\
\delta(q_1, \TMblank)  & = \tuple{q_1, \TMblank, \TMright}
\end{align*}
\end{ex}

\begin{explain}
Cette machine ne s'arrête que si l'entrée contient
un nombre pair de traits. Sinon, elle continue,
en théorie, de fonctionner indéfiniment. Pour
déterminer le résultat d'une machine sur une entrée,
on peut suivre la succession de ses \emph{configurations}.
Nous en donnerons plus loin une définition formelle.
Pour le moment, imaginons simplement une suite de
diagrammes décrivant la machine à chaque instant
du calcul. Une configuration indique le contenu
du ruban, l'état courant et la position de la
tête de lecture-écriture.

Suivons la machine paire sur une entrée de quatre
symboles $\TMstroke$. Nous nous attendons à ce
qu'elle s'arrête. Nous l'exécuterons ensuite sur
trois symboles $\TMstroke$ : elle fonctionnera
alors indéfiniment.

La machine commence dans l'état~$q_0$, avec la
tête sur le premier~$\TMstroke$. Sa configuration
initiale se note :
\[
\TMendtape \TMstroke_0 \TMstroke \TMstroke \TMstroke \TMblank \ldots
\]
L'indice~$0$ placé sur le premier~$\TMstroke$
indique que la tête lit ce symbole dans l'état~$q_0$.\footnote{Correction éditoriale : la glose source parle ici de l'état « one », mais le diagramme, la configuration et l'instruction indiquent tous $q_0$.}
L'instruction applicable est
$\delta(q_0,\TMstroke)=\tuple{q_1,\TMstroke,\TMright}$ :
la tête se déplace donc d'une case vers la droite
et la machine passe dans l'état~$q_1$.
\[
\TMendtape \TMstroke \TMstroke_1 \TMstroke \TMstroke \TMblank \ldots
\]
Comme la tête lit maintenant un~$\TMstroke$
dans l'état~$q_1$, elle suit l'instruction
$\delta(q_1,\TMstroke)=\tuple{q_0,
  \TMstroke,\TMright}$. On obtient :
\[
\TMendtape \TMstroke \TMstroke \TMstroke_0 \TMstroke \TMblank \ldots
\]
Les deux instructions s'appliquent de nouveau
dans le même ordre, ce qui donne les deux
configurations suivantes :
\[
\TMendtape \TMstroke \TMstroke \TMstroke \TMstroke_1 \TMblank \ldots
\]
\[
\TMendtape \TMstroke \TMstroke \TMstroke \TMstroke \TMblank_0 \ldots
\]
La machine lit alors un~$\TMblank$ dans
l'état~$q_0$. Le diagramme ne prévoit aucune
instruction pour cette situation : la machine
s'est arrêtée, et l'entrée est donc acceptée.

Partons maintenant d'une entrée de trois symboles
$\TMstroke$. Les premières instructions sont les
mêmes; seul le contenu du ruban change :
\[
\TMendtape \TMstroke_0 \TMstroke \TMstroke \TMblank \ldots
\]
\[
\TMendtape \TMstroke \TMstroke_1 \TMstroke \TMblank \ldots
\]
\[
\TMendtape \TMstroke \TMstroke \TMstroke_0 \TMblank \ldots
\]
\[
\TMendtape \TMstroke \TMstroke \TMstroke \TMblank_1 \ldots
\]
Après avoir franchi tous les $\TMstroke$,
la tête lit un~$\TMblank$ dans l'état~$q_1$.
Le diagramme prescrit alors
$\delta(q_1,\TMblank)=\tuple{q_1,
\TMblank,\TMright}$. Comme toutes les cases
suivantes du ruban contiennent $\TMblank$,
la machine répète cette instruction
\emph{indéfiniment}, reste dans l'état~$q_1$
et se déplace toujours plus à droite. Elle
ne s'arrête jamais et n'accepte donc pas l'entrée.
\end{explain}

\begin{explain}
Toutes les machines ne s'arrêtent pas. Si une
machine s'arrête quand aucune instruction n'est
applicable, on peut construire une machine qui
ne s'arrête jamais en prévoyant, dans chaque état,
une flèche pour chaque symbole. Pour que la machine
paire fonctionne indéfiniment même sur une entrée
paire, il suffit d'ajouter une instruction pour
la lecture de $\TMblank$ dans l'état~$q_0$.
\end{explain}

\begin{ex}
\[
\begin{tikzpicture}[->,>=stealth',shorten >=1pt,auto,node distance=2.8cm,
                    semithick]
  \tikzstyle{every state}=[fill=none,draw=black,text=black]

  \node[initial,state] (A)              {$q_0$};
  \node[state]         (B) [right of=A] {$q_1$};

  \path
  (A) edge [bend left] node {\TMtrans{\TMstroke}{\TMstroke}{\TMright}} (B)
      edge [loop above] node {\TMtrans{\TMblank}{\TMblank}{\TMright}} (A)
  (B) edge [loop above] node {\TMtrans{\TMblank}{\TMblank}{\TMright}} (B)
      edge [bend left] node {\TMtrans{\TMstroke}{\TMstroke}{\TMright}} (A);
\end{tikzpicture}
\]
\end{ex}

\begin{explain}
On peut également représenter une machine de
Turing par une \emph{table de transition}. Les
symboles de l'alphabet du ruban figurent en
colonnes et les états de la machine en lignes.
À l'intersection d'un état et d'un symbole, la
case donne le reste de l'instruction : nouvel
état, nouveau symbole et sens du déplacement.
La table permet de repérer aisément les situations
où la machine s'arrête : toute case vide marque
une absence d'instruction. Ces points d'arrêt ne
sont pas toujours aussi visibles dans un diagramme
d'états ou une liste d'instructions.
\end{explain}

\begin{ex}
Voici la table de la machine paire\footnote{Correction éditoriale : la source inverse le nouvel état et le symbole écrit dans les cases remplies de cette table. Nous suivons l'ordre annoncé dans le texte : nouvel état, symbole écrit, direction.} :
\[
\centering
\begin{tabular}{lllll}
\cline{2-4}
\multicolumn{1}{l|}{}      & \multicolumn{1}{c|}{$\TMblank$}
& \multicolumn{1}{c|}{$\TMstroke$}     & \multicolumn{1}{c|}{$\TMendtape$}   &  \\ \cline{1-4}
\multicolumn{1}{|l|}{$q_0$} & \multicolumn{1}{l|}{}
& \multicolumn{1}{l|}{\TMtrans{q_1}{\TMstroke}{\TMright}} &
\multicolumn{1}{l|}{\phantom{\TMtrans{q_1}{\TMstroke}{\TMright}}} &  \\ \cline{1-4}
\multicolumn{1}{|l|}{$q_1$} & \multicolumn{1}{l|}{\TMtrans{q_1}{\TMblank}{\TMright}}
& \multicolumn{1}{l|}{\TMtrans{q_0}{\TMstroke}{\TMright}} &
\multicolumn{1}{l|}{\phantom{\TMtrans{q_1}{\TMstroke}{\TMright}}} &  \\ \cline{1-4}
\end{tabular}
\]
On voit que la machine s'arrête lorsqu'elle
lit~$\TMblank$ dans l'état~$q_0$.
\end{ex}

\begin{explain}
Jusqu'ici, nous n'avons considéré que des machines
qui lisent et acceptent une entrée. Une machine de
Turing peut aussi écrire. Considérons, parmi
beaucoup d'autres exemples possibles, une
\emph{machine doubleuse} : si le ruban contient
initialement un bloc de~$n$ symboles $\TMstroke$,
elle produit un bloc de~$2n$ symboles $\TMstroke$.
\end{explain}

\begin{ex}
  \ollabel{ex:doubler} Avant de construire une machine
  doubleuse, il faut choisir une \emph{stratégie}.
  Notre machine ne peut pas mémoriser dans ses seuls états un nombre arbitrairement grand de
  symboles $\TMstroke$ déjà lus; elle doit donc
  les suivre sur le ruban. Une méthode consiste à
  séparer l'entrée de la sortie par un~$\TMblank$.
  La machine efface le premier~$\TMstroke$ de
  l'entrée, traverse le reste de celle-ci, laisse
  une case~$\TMblank$ puis écrit deux nouveaux
  symboles $\TMstroke$. Elle revient ensuite
  chercher le second~$\TMstroke$ de l'entrée
  et le double à son tour. Pour chaque symbole
  $\TMstroke$ de l'entrée, elle en écrit donc
  deux dans la sortie. L'effacement progressif
  de l'entrée garantit qu'aucun symbole n'est
  oublié ni traité deux fois. Une fois l'entrée
  entièrement effacée, le ruban contient un bloc
  de~$2n$ symboles $\TMstroke$. Le diagramme de
  la machine obtenue figure dans la \olref{fig:doubler}.
  \begin{figure}
\[
\begin{tikzpicture}[->,>=stealth',shorten >=1pt,auto,node distance=2.8cm,
                    semithick]
  \tikzstyle{every state}=[fill=none,draw=black,text=black]

  \node[initial,state] (1)              {$q_0$};
  \node[state]         (2) [right of=1] {$q_1$};
  \node[state]         (3) [right of=2] {$q_2$};
  \node[state]         (4) [below of=3] {$q_3$};
  \node[state]         (5) [left of=4]  {$q_4$};
  \node[state]         (6) [left of=5]  {$q_5$};

  \path (1) edge node {\TMtrans{\TMstroke}{\TMblank}{\TMright}} (2)
    (2) edge [loop above] node {\TMtrans{\TMstroke}{\TMstroke}{\TMright}} (2)
      edge node {\TMtrans{\TMblank}{\TMblank}{\TMright}} (3)
    (3) edge [loop above] node {\TMtrans{\TMstroke}{\TMstroke}{\TMright}} (3)
        edge node {\TMtrans{\TMblank}{\TMstroke}{\TMright}} (4)
    (4) edge [loop below] node {\TMtrans{\TMblank}{\TMstroke}{\TMleft}} (4)
        edge node {\TMtrans{\TMstroke}{\TMstroke}{\TMleft}} (5)
    (5) edge [loop below]  node {\TMtrans{\TMstroke}{\TMstroke}{\TMleft}} (5)
        edge              node {\TMtrans{\TMblank}{\TMblank}{\TMleft}} (6)
    (6) edge [loop below] node {\TMtrans{\TMstroke}{\TMstroke}{\TMleft}} (6)
        edge              node {\TMtrans{\TMblank}{\TMblank}{\TMright}} (1);
\end{tikzpicture}
\]
\caption{Une machine doubleuse}
\ollabel{fig:doubler}
\end{figure}
\end{ex}

\begin{prob}
Choisissez une entrée et suivez les configurations
successives de la machine doubleuse de l'\olref[tur][mac][rep]{ex:doubler}.
\end{prob}

\begin{prob}
Concevez une machine de Turing d'alphabet
$\{\TMendtape,\TMblank,A,B\}$ qui accepte,
c'est-à-dire s'arrête sur, toute chaîne de $A$
et de $B$ contenant autant de $A$ que de $B$
\emph{et} dans laquelle tous les $A$ précèdent
tous les $B$. Elle doit rejeter, c'est-à-dire
ne pas s'arrêter sur, les chaînes où ces nombres
diffèrent ou où un $A$ vient après un~$B$.
Par exemple, elle doit accepter $AABB$ et $AAABBB$,
mais rejeter $AAB$ et $AABBAABB$.
\end{prob}

\begin{prob}
Concevez une machine de Turing d'alphabet
$\{\TMendtape,\TMblank,A,B\}$ qui prend pour
entrée une chaîne quelconque~$\alpha$ de $A$
et de $B$, et la duplique pour produire
$\alpha\alpha$. Par exemple, l'entrée $ABBA$
doit produire $ABBAABBA$.
\end{prob}

\begin{prob}
\emph{Ordonnée ?} Concevez une machine de Turing
d'alphabet $\{\TMendtape,\TMblank,A,B\}$ qui,
sur une suite finie de $A$ et de $B$, détermine
si tous les $A$ précèdent tous les $B$. Elle
doit laisser la chaîne d'entrée sur le ruban,
puis s'arrêter si la chaîne est « ordonnée »
et tourner indéfiniment sinon.
\end{prob}

\begin{prob}
\emph{Trieuse} : concevez une machine de Turing
d'alphabet $\{\TMendtape,\TMblank,A,B\}$ qui
réordonne une suite finie de $A$ et de $B$ de
façon que tous les $A$ soient à gauche de tous
les~$B$. Par exemple, $BABAA$ doit devenir
$AAABB$, et $ABBABB$ doit devenir $AABBBB$.
\end{prob}

\end{document}
